\chapter{Introduction}\label{ch}

\section{Selecting Communication Stacks for Distributed GPU Applications}

Graphics Processing Units (GPUs) provide the parallelism and memory bandwidth used by scientific simulations and large-scale artificial-intelligence training and inference. As these workloads grow beyond one accelerator, they also require more communication: GPUs exchange gradients, activations, parameters, and routed tokens within a node and across cluster interconnects. Pilz et al.'s dataset of 500 AI supercomputers from 2019 to 2025 documents rapid growth in advertised computational capacity, hardware cost, and power demand \cite{pilz_trends_2025}. Those infrastructure trends motivate the scale of the problem; they do not by themselves establish an application communication bottleneck. At application level, performance depends on how data movement and synchronization interact with computation \cite{unat_landscape_2026,choi_end--end_2020}.

Exascale system throughput \cite{top500_june2026} does not remove this coordination problem. Even when a transfer accesses GPU memory directly, the CPU may still issue the operation, wait for a producing kernel, and arrange completion before the next kernel starts. Repeated host orchestration can leave devices idle or prevent useful overlap. Conversely, a host-driven library can remain efficient when its algorithms, batching, and progress mechanisms match the workload. The question is therefore when moving communication control into GPU kernels improves the critical path, rather than whether a system supports direct GPU data movement at all \cite{unat_landscape_2026,agostini_gpudirect_2018}.

This demand has accelerated the development of several programming models and communication libraries. MPI remains the predominant standard for distributed-memory programming and provides a portable interface for point-to-point and collective communication. GPU-aware MPI implementations additionally allow communication operations to access device memory directly, avoiding explicit staging through host memory.

Vendor-optimized collective libraries, such as NVIDIA NCCL and AMD RCCL, provide specialized implementations of operations such as all-reduce, broadcast, reduce-scatter, and all-gather. These libraries are particularly important in distributed deep-learning workloads, where collective operations frequently dominate communication time \cite{weingram_xccl_2023}. Their performance depends on several internal decisions, including the selected communication algorithm, transport protocol, number of communication channels, and underlying system topology \cite{hu_demystifying_2025}.

A different approach is provided by one-sided Partitioned Global Address Space models such as OpenSHMEM. GPU-oriented implementations, including NVSHMEM and ROC\_SHMEM, extend this model by allowing communication operations to be initiated directly by GPU threads or thread blocks. Device-initiated communication can reduce CPU involvement and enable fine-grained communication from inside a kernel, potentially improving overlap and reducing synchronization overhead \cite{agostini_gpudirect_2018,ma_demystifying_2026}. Portable implementations such as OSHMPI instead provide OpenSHMEM semantics over MPI. These alternatives make different trade-offs: symmetric-memory requirements, synchronization semantics, message granularity, host involvement, and implementation maturity can all determine whether one-sided communication is beneficial.

Recent work has demonstrated that default collective algorithms and runtime configurations can be significantly slower than configurations selected for a particular topology and workload \cite{pasqualoni_pico_2026,li_tuccl_2025}. The communication backend should therefore be treated as a workload-dependent design choice rather than as an interchangeable implementation detail.

The rapid growth of this ecosystem creates a practical selection problem. MPI, NCCL, oneCCL, NVSHMEM, and OSHMPI expose different interfaces and may use different algorithms, progress engines, memory requirements, and transports. Their names alone do not tell an application developer which stack will perform best. The answer depends on the application's communication signature: the operations it issues, their message sizes, peer sets, synchronization points, rank placement, and opportunities for computation--communication overlap.

Vendor-specific ecosystems generally provide highly optimized implementations, but they couple applications to a particular hardware and software stack. SYCL and oneCCL provide source-level abstractions that reduce this coupling. Source portability, runtime backend interchangeability, and performance portability are nevertheless distinct properties: a common interface may add orchestration, staging, or synchronization costs, and measurements on NVIDIA A100 GPUs cannot by themselves establish cross-vendor performance portability \cite{hammond_comparative_2019,shen_sycl_2026}. This thesis therefore treats CUDA and SYCL variants as concrete implementation stacks and supporting controls, rather than making cross-vendor portability a separate experimental claim.

\subsection{Research Gap and Approach}

Prior work already connects communication mechanisms to applications. Agostini et al. examine GPUDirect Async through microbenchmarks, modeling, and application cases; PICO assesses benchmark-informed choices through simulated training-trace replay; TuCCL reports end-to-end training improvements; and Trotter et al. compare host- and device-initiated CG solvers \cite{agostini_gpudirect_2018,pasqualoni_pico_2026,li_tuccl_2025,trotter_cpu-_2025}. The gap addressed here is the specific combination of a common CUDA/SYCL comparison across MPI, collective-library, and SHMEM paths with measured validation in an extended sparse solver. Supporting that comparison also requires functionality absent from the selected software baselines: oneCCL lacks the added OSHMPI and NVSHMEM paths and complete grouped NCCL dispatch, while the published aCG suite provides CUDA/HIP implementations without the OSHMPI communicator or SYCL port developed here.

This thesis addresses those gaps through a pattern-oriented workflow. It first compares configured host- and GPU-initiated paths to locate favorable regimes. It extends oneCCL with grouped NCCL support and an OSHMPI backend, and establishes a separately validated NVSHMEM runtime foundation. Finally, OSHMPI and SYCL extensions to aCG test the benchmark guidance in a complete sparse solver. The central hypothesis is conditional: persistent device work can amortize host-control costs, but memory placement, completion semantics, proxy progress, and application dependencies determine whether that opportunity becomes a performance benefit. Comparisons of complete stacks identify useful candidates; they do not isolate initiation as the only changing variable.

This is why the benchmark suite developed in this work is organized around communication patterns rather than around library interfaces. It covers reductions, neighbor exchanges, and all-to-all redistribution because they expose different requirements: reductions combine values globally, neighbor exchanges communicate with a small application-defined set of peers, and all-to-all redistributes distinct data among all ranks. Data-parallel and sharded training motivate the reduction family, Mixture-of-Experts expert routing motivates all-to-all and variable-count redistribution, and iterative sparse solvers motivate latency-sensitive reductions and halo exchanges. A pattern-oriented characterization can therefore inform several workload families without implying that every workload contains every pattern.

\section{Scope and Main Findings}

This thesis applies that workflow on the Leonardo Booster partition. It first measures point-to-point, halo, reduction, and all-to-all patterns across CUDA-aware and SYCL-aware MPI, NCCL, oneCCL, NVSHMEM, and OSHMPI stacks on 1 to 32 NVIDIA A100 GPUs. Most benchmark paths are host-initiated, while the NVSHMEM ping-pong and halo implementations issue communication from persistent cooperative GPU kernels. A composite \texttt{cg\_step} benchmark then combines a host-driven halo exchange with the small reductions characteristic of one Conjugate Gradient iteration. Finally, the study evaluates host-initiated and monolithic device-initiated aCG solvers on two large sparse matrices. Conjugate Gradient combines sparse matrix--vector multiplication, vector operations, halo exchanges, and global reductions; its low computational intensity and repeated synchronization make scalability sensitive to communication overhead \cite{trotter_cpu-_2025}. All-to-all remains a standalone pattern motivated by redistribution-heavy workloads such as Mixture-of-Experts and is not validated by aCG.

The main campaign produces a pattern-specific ranking. HPC-X MPI has the lowest point-to-point and isolated-halo startup cost; device-initiated NVSHMEM leads large steady-state halos; NCCL-based paths lead large allreduce and multi-node all-to-all. MPI/UCC prevents a 17--34$\times$ small-message all-to-all slowdown at 16 and 32 ranks, while its multi-node allreduce effect changes from a penalty below 64\,KiB to a benefit at and above that size. Direct OSHMPI with host-staged scalar reductions wins every multi-rank synthetic \texttt{cg\_step}; NCCL leads the communication-free controls. Chapter~\ref{ch:m-results} explains the completion, staging, and phase-timing qualifications behind these findings.

The predictive value of these benchmarks is tested with aCG. This work adds an OSHMPI communicator whose halo exchange uses OpenSHMEM puts while its scalar CG reductions retain CUDA-aware \texttt{MPI\_Allreduce}. The direct OSHMPI \texttt{cg\_step} instead uses OpenSHMEM reductions, so those reduction measurements have no matched aCG counterpart. For the matched paths, \texttt{cg\_step} predicts aCG's eight-byte reduction cost within 6.5--9.7\% for MPI and NCCL and within 14--21\% for NVSHMEM across a 16-fold range in process count. In contrast, the regular two-neighbor eager halo does not predict aCG's irregular many-neighbor rendezvous exchange, whose communication is largely hidden by sparse matrix--vector multiplication. NCCL remains the fastest stable host-initiated backend in aCG, while the MPI halo exhibits fast and slow execution regimes that also affect the subsequent allreduce.

A separate monolithic NVSHMEM aCG variant examines GPU-initiated execution and scales better than the blocking OSHMPI path at larger process counts. It also changes solver organization and is therefore not a strict backend substitution. Moreover, Leonardo does not enable IBGDA, so its inter-node RMA uses the \texttt{ibrc} CPU proxy and does not measure GPU-posted inter-node communication. The combined evidence locates the benefit more narrowly than a library ranking: GPU initiation helps persistent neighbor work and can improve solver scaling when it removes repeated host synchronization, but it does not overcome MPI's startup advantage and remains sensitive to proxy progress. Microbenchmarks predict application behavior only when communication primitives, peer structures, completion boundaries, and overlap opportunities are comparable. The rankings are specific to Leonardo; the decision method is reusable on other systems.

A later exploratory \texttt{nccl-tests} comparison examines NCCL~2.18.5 and 2.31.2 with ordinary and symmetrically registered buffers on four GPUs within one node. Its summary indicates promising small-allreduce improvements, with regressions at other sizes and little small-message all-to-all benefit. Section~\ref{sec:nccl-preliminary-results} reports this as a separate preliminary study: the available summary does not establish the selected kernel or application-visible initiation mode, and these measurements do not revise the main campaign's rankings.

\section{Research Questions}

The investigation follows the same progression as the contributions. It first establishes a common comparison across communication models, then asks what is required to expose missing paths through oneCCL, and finally tests the benchmark guidance after extending a complete solver. This ordering makes the novelty relative to the baseline software explicit:

\begin{researchquestion}[RQ1 -- Conditions for beneficial GPU initiation]
Under which communication patterns, message sizes, topologies, and completion regimes do the evaluated GPU-initiated paths become beneficial relative to host-initiated stacks, and what evidence supports launch, synchronization, or proxy-progress costs as explanations for the observed boundaries?
\end{researchquestion}

\begin{researchquestion}[RQ2 -- Extending oneCCL beyond its baseline backends]
Compared with the baseline oneCCL backend set, what functionality and semantic adaptations are required for grouped NCCL point-to-point, OSHMPI/OpenSHMEM execution, and an NVSHMEM backend foundation, and what asynchrony, staging, portability, and readiness trade-offs result?
\end{researchquestion}

\begin{researchquestion}[RQ3 -- Transfer to extended CG solvers]
After extending aCG with OSHMPI and SYCL implementations, to what extent do the benchmark's host-versus-GPU-initiation findings transfer to complete CG solvers, and how do irregular peers, synchronization, computation--communication overlap, memory placement, and solver organization explain agreement or failure?
\end{researchquestion}

\section{Major contributions}

\begin{keypoint}[Central proposition]
GPU initiation is a mechanism, not a performance guarantee. Persistent device work can amortize launch costs and reduce repeated host synchronization, while application overlap can hide communication. These opportunities pay off only when the resulting cost is lower than the competing path. Pattern-oriented microbenchmarks locate candidate regimes; a complete application checks whether peers, completion semantics, and dependencies preserve them.
\end{keypoint}

To answer these questions, this thesis makes five main contributions. The second and third extend oneCCL beyond the baseline backend behavior, while the fourth and fifth extend aCG beyond the published CUDA/HIP suite.

\begin{enumerate}

\item \textbf{A unified benchmark across vendor-specific and portable APIs.}
A set of representative communication patterns characterizes latency, bandwidth, synchronization overhead, message-size sensitivity, and topology-dependent behavior across seven runtime stacks. Unlike separate per-library tests, the harness applies common process-to-GPU mapping, warm-up, measurement, validation, and result-reduction policies to CUDA- and SYCL-based MPI, NCCL/oneCCL, NVSHMEM, and OSHMPI paths. Its NVSHMEM ping-pong and halo kernels provide the GPU-initiated cases, while host-initiated stacks establish the comparison surface. Global reductions and neighbor exchanges correspond directly to aCG phases and support RQ3. All-to-all and alltoallv remain standalone redistribution patterns motivated by deep-learning workloads. The design spans vendor-specific and portable interfaces, while the reported validation remains specific to NVIDIA A100 GPUs. The suite is available at \url{https://github.com/y3rbiadit0/gpu-comm-benchmark}.

\item \textbf{Grouped point-to-point support for the oneCCL NCCL backend.}
oneCCL exposes \texttt{ccl::group\_start} and \texttt{ccl::group\_end} so that a sequence of operations can be submitted as one batch. For the NCCL backend, grouping is required to submit mutually dependent point-to-point operations together and is also a performance mechanism: NCCL can execute the batch in one launch and assign grouped transfers to separate channels where possible \cite{hu_demystifying_2025}. Before this work the group lifecycle was implemented only for the native path. This contribution adds backend-aware dispatch to \texttt{ncclGroupStart} and \texttt{ncclGroupEnd}, nested-group handling, events that track the deferred NCCL work, explicit failure propagation, and group-aware send/receive fallbacks for all-to-all. The semantics are documented and covered by dedicated grouped point-to-point tests. The implementation is available at \url{https://github.com/y3rbiadit0/oneCCL/tree/fix/nccl_group_support}.

\item \textbf{SHMEM-family backend extensions for oneCCL.}
The complete OSHMPI path, selected with \texttt{CCL\_BACKEND=oshmpi}, adds backend discovery, communicator and rank validation, datatype and reduction mapping, collective implementations, symmetric staging for ordinary host or SYCL USM buffers, and put-based point-to-point operations with explicit completion. It exposes the practical cost of mapping oneCCL's asynchronous event interface onto blocking, host-driven OpenSHMEM operations and is measured in the benchmark. A separate \texttt{CCL\_BACKEND=nvshmem} prototype establishes optional build integration, CUDA-backed SYCL communicator selection, KVS-distributed UID bootstrap, ownership-safe lifecycle handling, and a fixed symmetric staging arena. Its runtime and staging foundation passed one- and two-node validation, but its oneCCL collective entry points remain explicitly unsupported and it contributes no performance data. The implementations are available at \url{https://github.com/y3rbiadit0/oneCCL/tree/feat/oshmpi} and \url{https://github.com/y3rbiadit0/oneCCL/tree/feat/nvshmem_collectives}.

\item \textbf{An OSHMPI communicator and application validation in aCG.}
The aCG solver is extended with an OSHMPI communicator to test whether benchmark-level one-sided halo results transfer to a complete sparse solver. The communicator uses OpenSHMEM for halo exchange while retaining CUDA-aware \texttt{MPI\_Allreduce} for scalar reductions, so the application comparison does not treat the direct OSHMPI \texttt{cg\_step} benchmark's OpenSHMEM reductions as a matched primitive. The solver is evaluated against its MPI, NCCL, host-NVSHMEM, and device-NVSHMEM variants. The complete campaign connects the benchmark patterns to sparse matrix--vector multiplication, overlap, synchronization, and convergence on two large SuiteSparse matrices. The communicator is available at \url{https://github.com/y3rbiadit0/aCG-oshmpi/tree/oshmpi}.

\item \textbf{A SYCL implementation of aCG.}
The CUDA-oriented application is additionally implemented through SYCL, with distributed execution through MPI and selectable MPI or oneCCL scalar reductions. This port provides a complete portable-compute control for determining whether kernel and orchestration costs can dominate a communication comparison. Its evaluation is kept separate from the CUDA communicator campaign and is used as supporting evidence rather than as a claim of cross-vendor performance portability. The implementation is available at \url{https://github.com/y3rbiadit0/acg-sycl}.

\end{enumerate}

\section{Thesis Organization}

Chapter~\ref{ch:m-problem} introduces the hardware, communication patterns, programming models, libraries, transports, and performance concepts required to interpret the measurements. Chapter~\ref{ch:m-methodology} maps research questions to artifacts and evidence, then describes the benchmark, aCG campaign, tuning, and preliminary NCCL protocol. Chapter~\ref{ch:m-results} reports microbenchmark, composite-step, and complete-solver evidence, followed by the separate exploratory NCCL comparison. Chapter~\ref{ch:m-related} positions the work relative to communication characterization, GPU initiation, and application-level studies. Chapter~\ref{ch:m-discussion} answers the research questions and states the threats to validity. The final chapter summarizes the conclusions and identifies the questions that require further experiments.

